\documentclass[10pt,a4paper]{article}
\usepackage{amsmath,amssymb,amsthm,graphicx}
\usepackage{amsmath,amssymb,amsthm,hyperref,xfrac}
\providecommand{\ie}{i.e.\ }
\providecommand{\skpt}{\smallskip}
\usepackage[scr=rsfs,cal=euler]{mathalfa}
\multlinegap0pt

\hyphenation{begin decom-po-si-tion sequence}
\newenvironment{enumeratea}
{\bgroup\def\theenumi{\alph{enumi}}\begin{enumerate}}
{\end{enumerate}\egroup}

\newcommand{\ba}{{\boldsymbol a}}
\newcommand{\bb}{{\boldsymbol b}}
\newcommand{\bd}{{\boldsymbol d}}
\newcommand{\bu}{{\boldsymbol u}}
\newcommand{\bx}{{\boldsymbol x}}
\newcommand{\bmy}{{\boldsymbol y}}

\newcommand{\bD}{{\boldsymbol D}}
\newcommand{\bL}{{\boldsymbol L}}
\newcommand{\bP}{{\boldsymbol P}}
\newcommand{\bV}{{\boldsymbol V}}
\newcommand{\bX}{{\boldsymbol X}}
\newcommand{\bZ}{{\boldsymbol Z}}

\newcommand{\Psfrac}[2]{(\sfrac{#1}{#2})}
\newcommand{\psfrac}[2]{\sfrac{(#1)}{#2}}
\newcommand{\bigO}[1]{O(#1)}

\usepackage{mathtools}
\usepackage{bbm}

\newcommand{\NN}{\mathbb{N}}
\newcommand{\ZZ}{\mathbb{Z}}
\newcommand{\RR}{\mathbb{R}}
\newcommand{\CC}{\mathbb{C}}
\newcommand{\PP}{\mathbb{P}}
\newcommand{\EE}{\mathbb{E}}

\newcommand{\boldalpha}{\boldsymbol{\alpha}}
\newcommand{\bolddelta}{\boldsymbol{\delta}}
\newcommand{\boldrho}{\boldsymbol{\rho}}
\newcommand{\dee}{\,\mathrm{d}}
\newcommand{\one}{\mathbbm{1}}
\newcommand{\BC}{\mathcal{B}}
\newcommand{\EC}{\mathcal{E}}
\newcommand{\RS}{\mathscr{R}}
\newcommand{\prob}[1]{\PP[ #1 ]}
\newcommand{\bigoo}[1]{O\!\left(#1 \right)}
\newcommand{\abs}[1]{\left\lvert #1 \right\rvert}
\newcommand{\floor}[1]{\left\lfloor #1 \right\rfloor}
\newcommand{\riemint}[4]{\int_{#1}^{#2}\! #3 \, \mathrm{d}#4}
\newcommand{\maxnorm}[1]{\left\| #1 \right\|_{\infty}}

\renewcommand{\Re}{\textup{Re}}
\renewcommand{\Im}{\textup{Im}}

\newcommand{\dtv}{\mathrm{d}_{\mathrm{TV}}}
\newcommand{\dw}{\mathrm{d}_{\mathrm W}}

\let\epsilon\varepsilon

\newcommand{\pextra}{P_{\mathtt{extra}}}
\newcommand{\parratia}{P_{\mathtt{Arratia}}}

\newcommand{\simplex}{\Delta^{k-1}}

\newcommand{\bg}{\big}
\newcommand{\bgg}{\Big}
\newcommand{\bggg}{\bigg}

\newcommand{\Rnt}{R_{\mathrm{NT}}}
\newcommand{\Rprob}{R_{\mathrm{PD}}}
\newcommand{\ECnt}{\EC_{\mathrm{NT}}}
\newcommand{\ECprob}{\EC_{\mathrm{PD}}}

\newcommand{\Bnt}{\mathcal{B}_{\mathrm{NT}}}
\newcommand{\Bprob}{\mathcal{B}_{\mathrm{PD}}}

\DeclareMathOperator{\vard}{Var}
\DeclareMathOperator{\Dir}{Dir}

\theoremstyle{plain}
\newtheorem{cor}{Corollary}[section]
\newtheorem{lem}[cor]{Lemma}
\newtheorem{prop}[cor]{Proposition}
\newtheorem{conj}[cor]{Conjecture}
\newtheorem{theo}{Theorem}

\theoremstyle{definition}
\newtheorem{example}{Example}
\newtheorem{deff}[cor]{Definition}
\newtheorem{rem}{Remark}[section]
\newtheorem*{rems}{Remark}
\newtheorem*{remss}{Remarks}

\makeatletter
\ams@newcommand{\multiint}[1]{\DOTSI\protect\MultiIntegral{#1}}
\renewcommand{\MultiIntegral}[1]{%
\edef\ints@c{\noexpand\intop
\ifnum#1=\z@\noexpand\intdots@\else\noexpand\intkern@\fi
\replicate{#1-2}{\noexpand\intop\noexpand\intkern@}%
\noexpand\intop
\noexpand\ilimits@
}%
\futurelet\@let@token\ints@a
}
\makeatother
\ExplSyntaxOn
\cs_new:Npn \replicate #1 #2 { \prg_replicate:nn { #1 } { #2 } }
\ExplSyntaxOff
\title{Uniform estimate for the smooth part of the multiplicative convolution sum}
\author{Tony Haddad, Dimitris Koukoulopoulos}
\date{doi:10.5802/jep.317; CC BY 4.0}
\begin{document}
\maketitle
\noindent\textit{Editorial note.} The counted statement is the author's unlabelled lem environment at source0.tex lines 1394-1399; its proof environment opens at line 1401, the line after the blank separator, so the pairing is by adjacency and the retained proof range 1401-1446 is the whole environment.\par\medskip
\noindent\textit{Editorial note.} Three environments are retained as context because the counted proof depends on them: the definition of the class of multiplicative functions with its labels (source lines 271-290), the description of the function on square-free arguments (source lines 294-297) and the partial summation bounds prel:partial summation bounds (source lines 1678-1685).\par\medskip
\noindent\textit{Editorial note.} Because the lemma carries no author label, its identity is recorded by the pinned source position of the two environments (statement lines 1394-1399, proof lines 1401-1446) and by the claim key, which encodes that position.\par\medskip
\noindent\textit{Editorial note.} This article carries no label-number dump and the retained counter values are left empty, so the printed numbering of the card is the wrapper's default; the original printed number of this unlabelled lemma is not recoverable from the pinned material and is therefore not asserted.\par\medskip
\noindent\textit{Editorial scope.} One original result of the pinned publisher article, counted once: the uniform estimate for the smooth part of the convolution sum (unlabelled lemma at source0.tex lines 1394-1399).\par\medskip
\noindent\textit{Allowed standard prerequisites (editorial).}\begin{itemize}
\item The class of multiplicative functions with a size condition. The functions of the class are multiplicative, and their values on square-free arguments are determined by the exponents attached to the prime factors, with a size condition on those exponents.. It is the class to which the dominated function belongs.
\item Square-full and square-free decomposition. Every integer factors uniquely as a product of its square-full part and a square-free part, and this decomposition is used to split the divisors.. It is the first step of the proof.
\item Partial summation. Sums over integers of a multiplicative function can be evaluated by partial summation once the corresponding summatory function is known.. It converts the estimate obtained for the sum into the stated bound.
\end{itemize}\medskip
\begin{deff}[The class of functions $\mathcal F_k(\boldalpha)$]\label{dfn:Fktheta}
Given $k\in\ZZ_{\ge2}$ and $\boldalpha\in \simplex$, let $\mathcal F_k(\boldalpha)$ be the set of functions $f \colon \NN^{k} \to \RR_{ \ge 0}$ satisfying the following three properties:
\begin{enumeratea}
\item\label{dfn:Fkthetaa} We have $\sum_{d_1\cdots d_k = n} f(d_1, \ldots, d_k) = 1$ for all $n \in \NN$. Thus, for each fixed $n \in \NN$, the function $f(\bd)$ defines a probability mass function on the set of vectors $\bd \in \NN^k$ satisfying $d_1\cdots d_k = n$.

\item\label{dfn:Fkthetab} Whenever $\bd$ satisfies
\[
d_i \coloneqq \begin{cases}
p &\text{if $i = j$}, \\
1 &\text{if $i \ne j$}.
\end{cases}
\]
for some $1 \le j \le k$ and prime $p$, then $f(\bd) = \alpha_j$.

\item\label{dfn:Fkthetac} The function $f$ is \emph{multiplicative}, \ie for any vectors $\ba, \bb \in \NN^k$ with $a_1\cdots a_k$ and $b_1 \cdots b_k$ being co-prime, we~have the property
\[
f(a_1b_1, \ldots, a_kb_k) = f(a_1, \ldots, a_k) \cdot f(b_1, \ldots, b_k).
\]
\end{enumeratea}
\end{deff}

\begin{equation}
\label{eq:f on squarefrees}
f(d_1,\dots,d_k) = \prod_{j=1}^k \alpha_j^{\omega(d_j)}.
\end{equation}

\begin{prop}
\label{prel:partial summation bounds}
Uniformly for $\alpha\in[0,1]$ and $x,y\ge2$, we~have
\[
\sum_{n \le x} \alpha^{\omega(n)} \ll x(\log x)^{\alpha-1}
\quad \text{and} \quad \sum_{x/y<n \le xy} \frac{ \alpha^{\omega(n)}}{n} \ll (\log y)(\log(xy))^{\alpha-1}.
\]
\end{prop}

\begin{lem}
There exists a universal constant $x_0\ge2$ such that if $x\ge x_0$ and $z\in[2,(\log x)^4]$, then we have the uniform estimate
\[
S(z) \ll \frac{\log z}{\sqrt{z}} \cdot \frac{x}{(1+ u\log x)^{1-\alpha}(1+(1-u)\log x)^{\alpha}}.
\]
\end{lem}

\begin{proof}
If $d_1\cdots d_k=n$, then we may uniquely write $d_j=d_j'd_j''$ with $d_j'|n^\flat$ and $d_j''|s(n)$. In~particular, $d_j''\le s(n)\le x/n^\flat\le z$ for all $j$. Hence, if $d_i\in[x^u z^{-3},x^uz^3]$, we~must also have that $d_i'\in[x^uz^{-4},x^uz^3]$. Hence, using property \eqref{dfn:Fkthetac} of Definition~\ref{dfn:Fktheta}, we~deduce that
\[
\sum_{\substack{d_1\cdots d_k = n \\ d_i\in[x^uz^{-3}, x^uz^3] }} \hspace*{-5mm}f(d_1,\dots,d_k)
\le \sum_{\substack{d_1'\cdots d_k' = n^\flat \\ d_i'\in[x^uz^{-4}, x^uz^3] }}\hspace*{-5mm} f(d_1',\dots,d_k')
\sum_{d_1''\cdots d_k'' = s(n)}\hspace*{-3mm} f(d_1'',\dots,d_k'').
\]
Relation \eqref{eq:f on squarefrees} implies that $f(d_1',\dots,d_k') = \prod_{j=1}^k \alpha_j^{\omega(d_j')}$. Moreover, using property \eqref{dfn:Fkthetaa} of Definition \ref{dfn:Fktheta}, we~find that the sum over the $d_j''$'s equals $1$. In~conclusion,
\begin{align*}
\sum_{\substack{d_1\cdots d_k = n \\ d_i\in[x^uz^{-3}, x^uz^3] }}\hspace*{-4mm} f(d_1,\dots,d_k)
&\le \sum_{\substack{d_1'\cdots d_k' = n^\flat \\ d_i'\in[x^uz^{-4}, x^uz^3] }} \prod_{j=1}^k \alpha_j^{\omega(d_j')} \\
&= \sum_{\substack{d| n^\flat \\ d\in[x^uz^{-4}, x^uz^3] }} \hspace*{-5mm} \alpha^{\omega(d)} (1-\alpha)^{\omega(n^\flat/d)}.
\end{align*}
Let $n^\flat=dm$ and $s(n)=b$. If~$n^\flat\in[\sfrac{x}{z},\sfrac{2x}{z}]$ and $d\in[x^uz^{-4},x^uz^3]$, then $m\in [x^{1-u}z^{-4},2x^{1-u}z^3]$. Hence, we~find that
\begin{equation}
\label{eq:bound in terms of S1 and S2}
\begin{aligned}
S(z) &\le \mathop{\sum\sum}_{\substack{dm \in[\sfrac{x}{z},\sfrac{2x}{z}] \\ d\in[x^uz^{-4}, x^uz^3] \\ m\in[x^{1-u}z^{-4},2x^{1-u}z^3] }}\hspace*{-5mm}
\alpha^{\omega(d)} (1-\alpha)^{\omega(m)} \sum_{\substack{b\le z \\ b\ \text{square-full}}}\hspace*{-5mm} 1\\
&\ll \sqrt{z}\hspace*{-5mm} \mathop{\sum\sum}_{\substack{dm \in[\sfrac{x}{z},\sfrac{2x}{z}] \\ d\in[x^uz^{-4}, x^uz^3] \\ m\in[x^{1-u}z^{-4},2x^{1-u}z^3] }}\hspace*{-4mm} \alpha^{\omega(d)} (1-\alpha)^{\omega(m)}
= \sqrt{z}\cdot (S_1+S_2),
\end{aligned}
\end{equation}
where $S_1$ denotes the double sum over $d$ and $m$ with the additional constraint \hbox{$d\le \sqrt{2x/z}$}, and $S_2$ is the corresponding sum over pairs $(d,m)$ with $d>\sqrt{2x/z}$ and $m\le \sqrt{2x/z}$.

First, we~bound $S_1$. Note that for the conditions on $d$ to be compatible we must have that $u\le 2/3$, provided that $x$ is large enough. Assuming this is the case, we~apply twice Proposition \ref{prel:partial summation bounds} to find that
\begin{equation}\label{eq:bound for S1}
\begin{aligned}
S_1 &\le \sum_{\substack{d\le \sqrt{2x/z} \\ d\in[x^uz^{-4}, x^uz^3] }}\hspace*{-5mm} \alpha^{\omega(d)} \sum_{m\le \sfrac{2x}{dz}} (1-\alpha)^{\omega(m)} \\
&\ll \sum_{\substack{d\le \sqrt{2x/z} \\ d\in[x^uz^{-4}, x^uz^3] }}\hspace*{-5mm} \alpha^{\omega(d)} \cdot \frac{x(\log x)^{-\alpha}}{dz} \\
&\ll \frac{x(\log z)(1+ u\log x)^{\alpha-1}(\log x)^{-\alpha}}{z}.
\end{aligned}
\end{equation}

We bound $S_2$ in a similar manner. For this sum to be non-empty, we~have $u\ge1/3$. If~this is the case, then we have
\begin{equation}
\label{eq:bound for S2}
\begin{aligned}
S_2
&\le \sum_{\substack{m\le \sqrt{2x/z} \\ m\in[x^{1-u}z^{-4}, 2x^{1-u}z^3] }}\hspace*{-5mm} (1-\alpha)^{\omega(m)} \sum_{d\le \sfrac{2x}{mz}} \alpha^{\omega(d)} \\
&\ll \sum_{\substack{m\le \sqrt{2x/z} \\ m\in[x^{1-u}z^{-4}, 2x^{1-u}z^3] }}\hspace*{-5mm} (1-\alpha)^{\omega(m)} \cdot \frac{x(\log x)^{\alpha-1}}{mz} \\
&\ll \frac{x(\log z)(1+ (1-u)\log x)^{-\alpha}(\log x)^{\alpha-1}}{z}.
\end{aligned}
\end{equation}
Combining \eqref{eq:bound in terms of S1 and S2}, \eqref{eq:bound for S1} and \eqref{eq:bound for S2}, while keeping in mind that $S_1=0$ if $u>2/3$ and that $S_2=0$ if $u<1/3$, completes the proof of the lemma,
\end{proof}

\begin{thebibliography}{99}
\bibitem{Arratia98} R. Arratia - “On the central role of scale invariant Poisson processes on (0,∞)”, in Microsurveys in discrete probability (Princeton, NJ, 1997), DIMACS Ser. Discrete Math. Theoret. Comput. Sci., vol. 41, American Mathematical Society, Providence, RI, 1998, p. 21-41
\bibitem{Arratia02} R. Arratia - “On the amount of dependence in the prime factorization of a uniform random integer”, in Contemporary combinatorics, Bolyai Soc. Math. Stud., vol. 10, János Bolyai Math. Soc., Budapest, 2002, p. 29-91
\bibitem{ABT06} R. Arratia, A. D. Barbour \& S. Tavaré - “A tale of three couplings: Poisson-Dirichlet and GEM approximations for random permutations”, Combin. Probab. Comput. 15 (2006) no. 1-2, p. 31-62
\bibitem{BareikisManstavicius07} G. Bareikis \& E. Manstavičius - “On the DDT theorem”, Acta Arith. 126 (2007) no. 2, p. 155-168
\bibitem{Bill72} P. Billingsley - “On the distribution of large prime divisors”, Period. Math. Hungar. 2 (1972), p. 283-289
\bibitem{delaBretecheTenenbaum16} R. de la Bretèche \& G. Tenenbaum - “Sur les processus arithmétiques liés aux diviseurs”, Adv. in Appl. Probab. 48 (2016) no. A, p. 63-76
\bibitem{DDT79} J.-M. Deshouillers, F. Dress \& G. Tenenbaum - “Lois de répartition des diviseurs. I”, Acta Arith. 34 (1979) no. 4, p. 273-285 (loose errata)
\bibitem{DonnellyTavare87} P. Donnelly \& S. Tavaré - “The population genealogy of the infinitely-many neutral alleles model”, J. Math. Biol. 25 (1987) no. 4, p. 381-391
\bibitem{Feng10} S. Feng - The Poisson-Dirichlet distribution and related topics. Models and asymptotic behaviors, Probability and its Appl., Springer, Heidelberg, 2010
\bibitem{JKB} N. L. Johnson, S. Kotz \& N. Balakrishnan - Discrete multivariate distributions, Wiley Series in Probab. and Stat.: Applied Probab. and Stat., John Wiley \& Sons, Inc., New York, 1997
\bibitem{Kingman93} J. F. C. Kingman - Poisson processes, Oxford Studies in Probab., vol. 3, The Clarendon Press, Oxford University Press, New York, 1993
\bibitem{KotzBal00} S. Kotz, N. Balakrishnan \& N. L. Johnson - Continuous multivariate distributions. Vol. 1. Models and applications, Wiley Series in Probab. and Stat.: Applied Probab. and Stat., Wiley-Interscience, New York, 2000
\bibitem{Koukoulo19} D. Koukoulopoulos - The distribution of prime numbers, Graduate Studies in Math., vol. 203, American Mathematical Society, Providence, RI, 2019
\bibitem{LastPenrose18} G. Last \& M. Penrose - Lectures on the Poisson process, Institute of Math. Stat. Textbooks, vol. 7, Cambridge University Press, Cambridge, 2018
\bibitem{Ken22} S.-K. Leung - “Dirichlet law for factorisation of integers, polynomials and permutations”, Math. Proc. Cambridge Philos. Soc. 175 (2023) no. 3, p. 649-676
\bibitem{KenWeb24} S.-K. Leung - “Dirichlet law with arbitrary parameters for factorizations of integers” (2025), Retrieved on 24 August 2025 from https://sites.google.com/view/sunkaileung/miscellaneous
\bibitem{MV07} H. L. Montgomery \& R. C. Vaughan - Multiplicative number theory. I. Classical theory, Cambridge Studies in Advanced Math., vol. 97, Cambridge University Press, Cambridge, 2007
\bibitem{NyandwiSmati13} S. Nyandwi \& A. Smati - “Distribution laws of pairs of divisors”, Integers 13 (2013), article ID A13, 13 pages
\bibitem{Tenenbaum00} G. Tenenbaum - “A rate estimate in Billingsley’s theorem for the size distribution of large prime factors”, Q. J. Math. 51 (2000) no. 3, p. 385-403
\end{thebibliography}
\end{document}
